\section{The History of Philosophy of Science as a Single Argument}

The history of philosophy of science is a single long argument about one
problem: Hume's problem of induction. Every major movement is either a
response to it, a refinement of a previous response, or an attempt to
dissolve it. The Bayesian/Kolmogorov framework is where the argument arrives.
We trace the through-line here, briefly; each movement clarifies what the
framework is solving.

\subsection{Hume: The Impossibility}

David Hume~\cite{hume1748} identified the deepest problem in the philosophy
of science. The problem of induction: no finite set of observations logically
justifies a universal law. The inference from ``all observed $F$s are $G$s''
to ``all $F$s are $G$s'' is not deductively valid. You can observe a million
white swans; the next one might be black. It was, in Australia.

Worse: the obvious response---induction is justified because it has worked in
the past---is itself an inductive argument. The justification is circular.
Hume concluded that inductive inference is a matter of psychological habit,
not reason.

\subsection{Popper: The Right Asymmetry, the Wrong Conclusion}

Karl Popper~\cite{popper1934} noticed that confirmation is logically weak
while falsification is logically strong. No number of white swans proves the
law; one black swan refutes it. He proposed falsifiability as the
demarcation criterion for science and characterized scientific progress as
conjecture and refutation: bold hypotheses, then honest attempts to falsify
them.

Popper was right about the asymmetry and right that unfalsifiable theories
are scientifically empty. He was wrong to conclude that science involves no
induction. His ``corroboration'' measure---how well a theory has survived
tests---is, precisely defined, the Bayesian likelihood ratio $P(E \mid H) /
P(E \mid \neg H)$: the update factor, stripped of the prior he refused to use.
Popper computed the right quantity and declined to multiply by the prior.
The Solomonoff prior completes his programme.

\subsection{Kuhn: The Right Phenomenology, No Normative Theory}

Thomas Kuhn~\cite{kuhn1962} argued that Popper's account does not match how
science actually works. Scientists operate within \emph{paradigms}---shared
frameworks of assumptions and methods---and do not abandon them at the first
anomaly. When anomalies accumulate to the point where confidence breaks down,
a revolution occurs: one paradigm replaces another, discontinuously.

Kuhn correctly described the phenomenology. He lacked a normative account:
when \emph{should} a paradigm shift occur? The MDL framework supplies the
answer (Section~\ref{sec:connections}).

\subsection{Lakatos: The Right Criterion, Incompletely Formalized}

Imre Lakatos~\cite{lakatos1978} synthesized Popper and Kuhn via the
methodology of scientific research programmes. Each programme has a
\emph{hard core} of protected assumptions and a \emph{protective belt} of
auxiliary hypotheses that absorb anomalies. A programme is \emph{progressive}
if its protective belt modifications generate new confirmed predictions;
\emph{degenerating} if they only explain known anomalies without predicting
new ones.

Under MDL, a progressive programme is one where successive modifications
decrease $L(H) + L(D \mid H)$ on new data. A degenerating programme is one
where modifications increase $L(H)$ without decreasing $L(D \mid H)$ on new
observations. Lakatos gave the right criterion; MDL makes it computable.

\subsection{Feyerabend: Methodological Anarchism as the Solomonoff Prior}

Paul Feyerabend~\cite{feyerabend1975} argued that no hypothesis should be
ruled out \emph{a priori}: science has always advanced by entertaining ideas
that violated current methodological rules. ``Anything goes'' as a prior
policy.

The Solomonoff prior assigns nonzero probability to every computable
hypothesis. Nothing is excluded before the data speaks. Feyerabend's
permissiveness is not anarchism---it is the correct prior policy. A prior
that excludes any computable hypothesis is provably suboptimal. What
Feyerabend called anarchism, Solomonoff called universality.

\subsection{Bayesian Confirmation: The Culmination}

Bayesian confirmation theory~\cite{howson2006} provides the most complete
response to Hume. Evidence $E$ confirms $H$ if and only if $P(H \mid E) >
P(H)$; the degree of confirmation is the likelihood ratio $P(E \mid H) /
P(E)$. This dissolves the raven paradox (observing a red fire engine confirms
``all ravens are black,'' but only infinitesimally), handles old evidence via
counterfactual reasoning, and validates Lakatos's distinction as a distinction
between increasing and decreasing marginal likelihood.

Carnap spent his career trying to construct an inductive logic---a formal
system that would determine the degree of confirmation of any hypothesis by
any evidence. He never succeeded because he had no principled account of the
prior. The Solomonoff prior completes the programme he started.